\chapter{Real-World Performance}
\label{chap:ch8}

\par Performance in Gentlix has two faces: how much code each backend needs to maintain 
(the development cost) and how fast it responds under realistic workloads (execution 
cost). Safety tooling and library choices from earlier chapters are not free; this chapter 
measures whether the Rust backend pays a visible price on the same machine and schema 
as the C backend.

\section{Development Performance}
\label{sec:ch8sec1}

\par Both backends implement the same product surface, yet the C tree is larger and carries 
more manual verification work. This section reports lines of code by layer, a Phase~1 
repository audit of safety obligations, and how each language's build system affects 
reproducibility and undefined-behaviour risk.

\subsection{Source size by architectural layer}
\label{subsec:ch8sec1loc}

\par Table~\ref{tab:loc_by_layer} groups source lines by architectural layer (excluding 
build artefacts). Domain logic is similar in size because business rules align; the gap 
widens in repository, service, and HTTP layers where C hand-rolls libpq mapping, JSON, 
JWT, and dashboard aggregation.

\begin{table}[htbp]
\centering
\begin{tabular}{|l|r|r|r|}
\hline
\textbf{Layer} & \textbf{C (lines)} & \textbf{Rust (lines)} & \textbf{$\Delta$} \\
\hline
Domain & 983 & 737 & +246 \\
Repository / DB & 1\,494 & 871 & +623 \\
Service & 3\,651 & 1\,996 & +1\,655 \\
HTTP / server & 2\,057 & 1\,477 & +580 \\
\hline
\textbf{Total} & \textbf{8\,185} & \textbf{5\,081} & \textbf{+3\,104} \\
\hline
\end{tabular}
\caption{Source lines per layer in Gentlix (Phase~1 audit)}
\label{tab:loc_by_layer}
\end{table}

\par The C \texttt{dashboard\_service} and analytics sections in \texttt{transaction\_service.c} 
account for much of the service-layer gap---parallel aggregation, hand-built maps, and JSON 
serialisation that Rust expresses with fewer lines plus standard collections.

\subsection{Repository audit: manual safety obligations}
\label{subsec:ch8sec1audit}

\par To quantify review load, both trees were scanned for patterns C must verify manually 
but safe Rust handles through types or crates (excluding build artefacts). Table~\ref{tab:audit_obligations} 
should be read as \emph{verification effort}, not typing speed.

\begin{table}[htbp]
\centering
\begin{tabular}{|l|r|r|}
\hline
\textbf{Metric} & \textbf{C backend} & \textbf{Rust backend} \\
\hline
Source files (.c/.h vs .rs) & 78 & 39 \\
Lines of source & 8\,185 & 5\,081 \\
\texttt{malloc} call sites & 19 & 0 (via \texttt{Box}/\texttt{Vec}) \\
\texttt{free} call sites & 152 & 0 (\texttt{Drop}) \\
Explicit NULL checks / uses & 198 & 0 (\texttt{Option}) \\
\texttt{unsafe} blocks & n/a & 0 in audited modules \\
Hand-rolled JWT (lines) & $\sim$294 & $\sim$62 (crate + middleware) \\
\texttt{PQgetvalue} column reads & 29 & 0 (\texttt{sqlx} macros) \\
Compile-time SQL macros & 0 & 29 \\
\hline
\end{tabular}
\caption{Manual safety obligations in Gentlix (Phase~1 repository audit)}
\label{tab:audit_obligations}
\end{table}

\par Every schema or API change in C must re-check those sites; Rust more often fails the 
build when a query or handler signature drifts (Chapter~\ref{subsec:ch6sec1-2}). Representative 
defect injections on real Gentlix paths---statement replay bounds, swapped identifiers, 
wrong column indices---compiled in C and could yield HTTP~200 with wrong financial data; 
the same changes in Rust typically failed to compile or panicked in a controlled way 
(Blackburn~\cite{Blackburn2025}).

\subsection{Build systems and reproducibility}
\label{subsec:ch8sec1build}

\par \textbf{C / CMake.} The C backend links \texttt{libpq}, libmicrohttpd, jansson, OpenSSL, 
and Argon2 through \texttt{find\_library} in \texttt{CMakeLists.txt}. Missing libraries fail 
at link time; wrong versions may only appear at runtime on a new machine. Header guards and 
include paths enforce layer boundaries by convention, not by the build graph.

\begin{lstlisting}[language=make, caption={CMakeLists.txt (excerpt) --- Each dependency
must be discovered on the host or in the Docker build stage.}]
find_library(LIBMICROHTTPD NAMES microhttpd REQUIRED)
find_library(LIBSSL NAMES ssl REQUIRED)
find_library(LIBCRYPTO NAMES crypto REQUIRED)
find_library(LIBARGON2 NAMES argon2 REQUIRED)
\end{lstlisting}

\par Gatchell~\cite{Gatchell2003} reminds C teams that build scripts and link lines are part 
of the safety story---a backend that compiles against the wrong OpenSSL header set can pass 
tests yet fail JWT verification in production.

\par \textbf{Rust / Cargo.} Dependencies live in \texttt{Cargo.toml}; \texttt{Cargo.lock} 
pins versions for reproducible CI and Docker builds. \texttt{cargo build} fetches crates, 
compiles, and links without a separate package manager step on the developer machine. 
McNamara~\cite{McNamara2021} notes that Cargo's unified workflow is one reason Rust services 
start small quickly; Gentlix still spent time aligning JSON field names with C, but rarely 
on link errors for missing JWT libraries.

\par Module folders (\texttt{domain/}, \texttt{db/}, \texttt{service/}, \texttt{server/}) 
mirror the C header boundaries from Chapter~\ref{sec:ch4sec5}. \texttt{sqlx} offline data 
in \texttt{.sqlx/} lets builds run without a live PostgreSQL instance---trading setup 
complexity for compile-time schema checks described in Section~\ref{subsec:ch6sec1-2}.

\subsection{Undefined behaviour versus defined panics}
\label{subsec:ch8sec1ub}

\par In C, integer overflow, out-of-bounds access, null dereference, and data races are 
undefined behaviour; optimisers may assume they never happen. In a ledger, silent wraparound 
or torn reads are worse than a crash. Gentlix uses explicit guards on \texttt{int64\_t} 
arithmetic in C and \texttt{checked\_mul} / validation in Rust for sensitive paths 
(Section~\ref{sec:ch5sec1}).

\par Safe Rust defines panics for bounds failures and rejects data races at compile time 
when sharing rules are violated (Chapter~\ref{sec:ch5sec3}). Drysdale~\cite{Drysdale2024} 
treats \texttt{unsafe} as an explicit audit surface; Gentlix's audited application modules 
contain no \texttt{unsafe} blocks, so memory and aliasing rules stay in the checked subset 
during development and benchmarks.

\subsection{Developer experience in practice}
\label{subsec:ch8sec1experience}

\par Qualitatively, the C backend required more custom infrastructure before feature work 
could start (\texttt{Db}, JWT, maps, mutex analytics). Rust reached feature parity faster 
after the borrow-checker learning curve but needed repeated passes to match C JSON exactly. 
Palmieri~\cite{Palmieri2022} describes that curve as front-loaded: once types encode the 
domain, changes become safer but not always faster to write than editing C strings.

\section{Execution Performance}
\label{sec:ch8sec2}

\subsection{Methodology}
\label{subsec:ch8sec2method}

\par Execution benchmarks use \textbf{k6} on a local Windows host against 
\texttt{http://localhost:6767}. PostgreSQL runs in Docker with migration scripts 
\texttt{001}--\texttt{003} and the benchmark seed (\(\sim\)100\,000 transactions for user 
\texttt{gentlix@benchmark.com}). Only \textbf{one backend runs at a time} on port~6767; 
each workload runs three times; reported values are the \textbf{median} of those runs.

\par Each script uses \textbf{15~s warmup} and \textbf{30~s steady} load. Workloads mirror 
real product paths:

\begin{itemize}
  \item \textbf{W1 (login)} --- 32 virtual users, \texttt{POST /api/users/login} (Argon2id);
  \item \textbf{W2 (statement)} --- 8 VUs, \texttt{GET /api/transactions/statement} with date filter;
  \item \textbf{W3 (analytics)} --- 8 VUs, \texttt{GET /api/transactions/summary/monthly?per\_account\_limit=200000};
  \item \textbf{W4 (transfer)} --- 16 VUs, \texttt{POST /api/transactions/transfer} (1 cent between owned accounts; database re-seeded between run batches).
\end{itemize}

\par Scripts and result summarisation live in \texttt{backend/benchmark/} in the repository. 
AWS deployment (Chapter~\ref{sec:ch2sec4}) is the production target; cloud benchmarks were 
not run for this thesis but could reuse the same k6 scripts against a load balancer URL.

\subsection{Measurement caveat}
\label{subsec:ch8sec2caveat}

\par During measurement the C backend ran as an optimised release-style binary; the Rust 
backend was started with \texttt{cargo run} (debug profile by default). That asymmetry 
favours C on CPU-heavy workloads and should be kept in mind when reading Table~\ref{tab:k6_results}. 
A follow-up with \texttt{cargo run --release} on AWS would narrow the build-mode gap.

\subsection{Results}
\label{subsec:ch8sec2results}

\begin{table}[htbp]
\centering
\begin{tabular}{|l|c|r|r|r|r|}
\hline
\textbf{Wkld} & \textbf{Conc.} & \textbf{C p95 (ms)} & \textbf{Rust p95 (ms)} &
\textbf{C rps} & \textbf{Rust rps} \\
\hline
W1 Login & 32 & 2\,013 & 3\,839 & 25.9 & 14.5 \\
W2 Statement & 8 & 1\,928 & 3\,495 & 5.0 & 3.2 \\
W3 Analytics & 8 & 2\,699 & 7\,310 & 3.3 & 1.6 \\
W4 Transfer & 16 & 247 & 1\,089 & 115.9 & 69.1 \\
\hline
\end{tabular}
\caption{Median k6 results on local Windows host (shared PostgreSQL, one backend at a time)}
\label{tab:k6_results}
\end{table}

\par On this setup \textbf{C leads on both tail latency (p95) and throughput (rps)} for 
every workload. W1 highlights CPU-bound login on connection threads versus Rust's blocking 
pool and async stack overhead (Section~\ref{subsec:ch6sec2-4}). W2 and W3 stress 
statement replay and parallel analytics from Chapter~\ref{sec:ch5sec3}; Rust's 
\texttt{spawn\_blocking} and \texttt{sqlx} async path add cost even when safety improves. 
W4 is the lightest database write path; C's lower p95 fits a small transfer with minimal 
allocation work in release mode.

\par These numbers do \emph{not} mean C is the better language for new FinTech backends in 
general. They show that, \textbf{for this prototype on this hardware}, the safety and 
abstraction choices in Rust had a measurable runtime cost, while the C backend carried 
a larger manual verification load (Tables~\ref{tab:audit_obligations} 
and~\ref{tab:loc_by_layer}). Production decisions would combine both tables with team 
experience, tooling, and deployment environment---including a fair release-vs-release 
comparison on AWS if time allows.

\section*{Summary}

\par Gentlix trades manual safety effort for raw speed in C, and compiler-checked structure 
for higher local latency in Rust on the measured workloads. Development metrics and k6 
results should be read together: the language that reduces audit burden is not automatically 
the language that wins every benchmark on a developer laptop.
